\textbf{Hand-Designed Agent Systems} \quad
Researchers have designed numerous agent systems tailored to various tasks based on predefined heuristics and prior knowledge. These systems often employ techniques such as prompt engineering~\citep{chen2023unleashing,schulhoff2024prompt}, chain-of-thought reasoning and planning~\citep{wei2022chain,yao2022react}, as well as reflection~\citep{shinn2024reflexion,madaan2024self}, code generation~\citep{wang2023voyager,vemprala2024chatgpt}, tool use~\citep{nakano2021webgpt,qu2024tool}, retrieval-augmented generation~\citep{lewis2020retrieval,zhang2024survey}, and multi-agent collaboration~\citep{xu2023expertpromptinginstructinglargelanguage,wu2023autogen,qian2023communicative,hong2023metagpt}. Once crafted by human designers, these systems remain static and do not adapt or evolve over time.

\noindent\textbf{Meta-Learning Optimized Agent Systems} \quad
Some researchers have explored methods for enhancing agents through fixed learning algorithms~\citep{zhou2024symbolic,hu2024automated}. For example, certain frameworks store an agent’s successful or failed strategies in memory based on environmental feedback~\citep{liu2023think,hu2023chatdb,qian2024investigateconsolidateexploitgeneralstrategyintertask}, while others automatically optimize agent prompts~\citep{khattab2023dspy,zhang2024agent,khattab2023dspy}. Some studies focus on designing prompts that enable agents to autonomously refine specific functions~\citep{zhangoffline}. 
% \citet{zhou2024symbolic} proposed a symbolic learning framework that uses natural language gradients to optimize the structure of agents. \citet{hu2024automated} used a basic meta agent to design agents for downstream tasks. 
However, these meta-algorithms are also designed manually and remain unchanged once deployed, limiting the agents' ability.

\noindent\textbf{Recursive Self-Improvement} \quad
The concept of recursive self-improvement has a long history~\citep{good1966speculations,schmidhuber1987evolutionary}.  \godel machine~\citep{schmidhuber2003godel} introduced the notion of a proof searcher that executes a self-modification, thereby enabling the machine to enhance itself.
%only if it can prove that the modification is optimal, thereby enabling the machine to enhance itself continuously. Subsequent works by \citet{nivel2013bounded} and \citet{steunebrink2016growing} proposed restrictive modifications to ensure safety during the self-improvement process. 
In the early days, there were also some discussions of self-improving agents that were not based on LLM~\citep{hall2007self,steunebrink2012towards}.
More recently, \citet{zelikman2023self} applied recursive self-improvement to code generation, where the target of improvement was the optimizer itself.
% , and the utility was evaluated based on performance in downstream tasks.
% Glore \citep{havrilla2024glorewhenwhereimprove} proposes Stepwise ORMs to improve LLM reasoning through global and local refinements. V-star \citep{hosseini2024vstartrainingverifiersselftaught} trains a verifier to evaluate both correct and incorrect self-generated solutions. 
% RISE \citep{qu2024recursiveintrospectionteachinglanguage} enables recursive self-improvement by fine-tuning models to introspect and correct previous mistakes in multiple iterations. 
% SCoRe \citep{kumar2024traininglanguagemodelsselfcorrect} uses reinforcement learning to improve self-correction in LLMs by learning from self-generated correction traces. 
Some work~\citep{havrilla2024glorewhenwhereimprove,qu2024recursiveintrospectionteachinglanguage,kumar2024traininglanguagemodelsselfcorrect} also explores recursive self-improvement by fine-tuning models to introspect and correct previous mistakes.
\framework represents the first self-referential agent based on LLM. This approach is more flexible, removing human-designed constraints.
% and allowing the agent’s capabilities to be limited only by the foundational model itself, rather than by human design bottlenecks.